% Minimal preamble for P2 cards taken from sources/comparability-partitions-2408.13523v2/source0.tex
% (arXiv:2408.13523v2, "Partitioning perfect graphs into comparability graphs").
% Every line below is copied mechanically from that file; nothing was invented, renamed or redefined.
%   line 2    : \documentclass[12pt]{article}
%   line 3    : \usepackage{amssymb,latexsym}
%   line 4    : \usepackage{amsmath}
%   line 5    : \usepackage[dvipdfm]{graphicx}  -- the dvipdfm driver option is dropped, because the
%               assembler forces the xelatex engine and a forced driver option contradicts it.
%   line 16   : \usepackage{amsfonts}
%   line 17   : \usepackage{amsthm}
%   lines 20-30: the author's theorem declarations
%   lines 32-41: the author's own macro definitions
% Omitted source preamble lines and why they are not needed for these body spans:
%   line 1 (provenance comment), lines 6-11 (tikz and its libraries: no drawing in the retained spans),
%   lines 7 and 9 (verbatim), lines 12-13 (page geometry), line 14 (rawfonts),
%   line 15 (psfrag), line 18 (hyperref with the pagebackref option: the retained spans print no
%   citation anchors and the assembly adds its own hyperref machinery), lines 43-61 (the author's
%   backref neutralisation and his margin-note macros, which depend on the omitted hyperref).
% The retained spans of the body use only \em, \bf, \cal, amsthm environments, amsmath displays and
% the macros defined in lines 32-41 below.
\documentclass[12pt]{article}
\usepackage{amssymb,latexsym}
\usepackage{amsmath}
\usepackage{graphicx}
\usepackage{amsfonts}
\usepackage{amsthm}

\newtheorem{theorem}{Theorem}%[section]
\newtheorem{problem}[theorem]{Problem}%[section]
\newtheorem{observation}[theorem]{Observation}
\newtheorem{lemma}[theorem]{Lemma}
\newtheorem{corollary}[theorem]{Corollary}
\newtheorem{definition}{Definition}
\newtheorem{example}[theorem]{Example}
\newtheorem{conjecture}{Conjecture}%[section]
\newtheorem{proposition}[theorem]{Proposition}
\newtheorem{claim}[theorem]{Claim}
\newtheorem{remark}[theorem]{Remark}

\def\qed{\ifhmode\unskip\nobreak\fi\quad\ifmmode\Box\else$\Box$\fi}
\def\pont{\hspace{-6pt}{\bf.\ }}
\def\qed{\hfill $\Box$}

\def\R{{\mbox{\boldmath$\mathcal R$}}}
% the author's second definition of \R, left commented out exactly as in the source
\def\F{{\mathcal F}}
\def\P{{\mathcal P}}

\def\newpage{\vfill\eject}
\title{A logarithmic lower bound for the comparability partition number of the interval graph G\_n}
\author{András Gyárfás, Márton Marits, Géza Tóth}
\date{arxiv:2408.13523v2; CC BY 4.0}
\begin{document}
\maketitle
\noindent\textit{Editorial note.} The counted unit is the article's own Theorem loglog at source lines 432-433 together with the author's proof that opens with the bold header Proof. at source line 438 and closes with the square at source line 453; the proof is preceded by a medskip, so the pairing is by adjacency through blank lines only.\par\medskip
\noindent\textit{Editorial note.} The retained context at source lines 128-129, 391-395 and 397-400 is the article's own definitions of the parameters c(G) and p(G), of the interval graph G\_n with the abbreviations c(n) and p(n), and of the nested, crossing and touching pairs of intervals.\par\medskip
\noindent\textit{Editorial note.} The retained context at source lines 421, 424-425 and 427-428 is the article's own introduction of the double shift graph DS\_n, its vertices and adjacency rule, and the corollary that its chromatic number is at least log log n; the intervening sentence of source line 426, which carries a bibliographic citation to Erdos and Hajnal, is not retained, so the material carries no bibliography.\par\medskip
\noindent\textit{Editorial note.} The source preamble cannot be retained as it stands: it loads rawfonts, psfrag, verbatim, tikz with several libraries, hyperref with the pagebackref option and a dvi driver option for graphicx, none of which the retained body spans use, and it redefines backref commands that only exist because of that hyperref option. The retained preamble file therefore carries, line for line, source line 2, source line 3, source line 4, source line 16, source line 17, the theorem declarations of source lines 20-30 and the author macro definitions of source lines 32-41; the only textual change anywhere is the removal of the dvipdfm driver option from the graphicx line, because the assembly forces the xelatex engine. No new macro, package or declaration was introduced.\par\medskip
\noindent\textit{Editorial note.} The article's theorems co-triangulated, unit, gensplit, bipline and ub, its lemma rec, its corollaries doblesh, shift and p', and its proposition H9 are not counted here; of these only the corollary on the chromatic number of DS\_n is retained, and it is retained as context for the quoted lower bound rather than as a premise proved in this material.\par\medskip
\noindent\textit{Allowed standard prerequisites (editorial).}\begin{itemize}
\item The cover number c(G) of a graph by comparability graphs. c(G) is the smallest number of comparability subgraphs of G needed to cover the edge set of G, without requiring them to be edge-disjoint; the article's own definition is retained at source line 129, together with p(G) at source line 128.. The counted theorem is a lower bound for exactly this parameter, so the object being bounded must be fixed before the statement can be read.
\item The interval graph G\_n and the numbers c(n) and p(n). For n at least 2, G\_n is the intersection graph of the n choose 2 closed intervals whose endpoints are among n consecutive positive integers, and the article abbreviates c(n)=c(G\_n) and p(n)=p(G\_n); retained at source lines 391-395.. The counted lower bound is the statement c(n) >= log log(n) / 2 for these graphs, and the proof works with a cover of G\_n by comparability subgraphs.
\item Nested, crossing and touching pairs of intervals. Two distinct closed intersecting intervals with integer endpoints are nested when one contains the other; otherwise they are crossing when their intersection has more than one element and touching when the intersection is a single point; retained at source lines 397-400.. The proof labels the touching edges of G\_n, so the notion of a touching pair and the shape of its endpoints i<j<k must be available.
\item The double shift graph DS\_n and its lower bound. The double shift graph has the triples (i,j,k) as vertices, with two triples joined when the last two entries of one are the first two entries of the other, and the article records at source lines 427-428 the corollary that its chromatic number is at least log log(n), obtained from a result of Erdos and Hajnal; the vertices and adjacency rule are retained at source lines 421 and 424-425.. The proof converts a cover of G\_n by t comparability subgraphs into a proper colouring of DS\_n by 2t colours, so both the graph and the quoted lower bound on its chromatic number must be available.
\end{itemize}\medskip
Let $p(G)$ denote the minimum number $m$ such that $E(G)=\cup_{i=1}^m E(G_i)$ where the $G_i$-s are edge-disjoint comparability subgraphs of $G$.
Dropping  the condition that the $G_i$-s are edge-disjoint, we define $c(G)$ as the smallest number of comparability subgraphs of $G$ needed to {\em cover} $E(G)$.

For $n\ge 2$, let $G_n$ be the interval graph defined as the intersection graph
of the $n\choose 2$ closed intervals with endpoints among $n$ consecutive positive integers,
called {\em base points}. Note that $|V(G_n)|={n\choose 2}$.
For convenience, set $c(n)=c(G_n),p(n)=p(G_n)$.
Since $G_2$ has one vertex and no edge, we set $c(2)=p(2)=0$.

Two distinct closed intersecting intervals $I_1,I_2$ with
integer endpoints are called {\em nested } if one contains the other.
If $I_1,I_2$ are non-nested then they are either {\em crossing}
(if  $|I_1\cap I_2|\ne 1$) or {\em touching} (if  $|I_1\cap I_2|=1$).

The {\em double shift graph}, $DS_n$, is a close relative of shift graphs,

Two vertices, $(i, j, k)$ and $(x, y, z)$ are connected by an edge if and only if
(i) $j=x$ and $k=y$ or (ii) $i=y$ and $j=z$. From a result of Erd\H os and Hajnal

\begin{corollary}\label{doblesh} $\chi(DS_n)\ge\log\log(n)$.
\end{corollary}

\begin{theorem}\label{loglog} $c(n)\ge  {1\over 2}\log\log(n)$.
\end{theorem}

\noindent{\bf Proof. }  Assume that $G_n$ is the union of $t$ comparability subgraphs,
$H_1,\dots,H_t$. Consider each $H_i$ with a transitive orientation.
     Assign a type to each touching edge $e=\{[i,j],[j,k]\}$
    as a pair $(m,+)$ or $(m,-)$ where the index $m$ indicates that the touching edge is
    in $H_m$ and the + or - indicates whether the touching edge is oriented from $[i,j]$
    to $[j,k]$ or backwards.
    (Since $H_1,\dots,H_t$ is a cover, we might have several choices for $m$, in this case
    take an arbitrary one.) The number of types is clearly $2t$.

    Note that the types of the touching edges $e=\{[i,j],[j,k]\}$ and $f=\{[j,k],[k,l]\}$ must be different
    for all $1\le i<j<k<l\le n$, otherwise the transitivity assumption implies that there is an edge of $G_n$
    from $[i,j]$ to $[k,l]$ - a contradiction since those intervals are disjoint.

    Therefore, we obtained a proper coloring of the double shift graph $DS_n$ by $2t$ colors:
    the color of the vertex $(i, j, k)$ corresponds to the type of
    the touching edge $e=\{[i,j],[j,k]\}$. Consequently, $2t\ge \chi(DS_n)\ge \log\log(n)$, finishing the proof.  \qed

\end{document}
